\documentclass[11pt,a4paper]{article}

\usepackage[utf8]{inputenc}
\usepackage[margin=1in]{geometry}
\usepackage{amsmath,amssymb}
\usepackage{enumitem}
\usepackage{fancyhdr}
\usepackage{xcolor}

% Header
\pagestyle{fancy}
\fancyhf{}
\lhead{\textbf{Modern Computer Vision} --- Technion, Spring 2026}
\rhead{HW 1 --- Theory}
\renewcommand{\headrulewidth}{0.4pt}

% Operators
\DeclareMathOperator{\conv}{*}
\DeclareMathOperator{\supp}{supp}

\begin{document}

\begin{center}
  {\LARGE\bfseries Homework 1 --- Theory}\\[6pt]
  \textcolor{gray}{Due: April 26, 2026}
\end{center}

\vspace{6pt}

\noindent
Recall the definition of continuous convolution:
\[
(f \conv g)(t) \;=\; \int_{-\infty}^{\infty} f(\tau)\, g(t - \tau)\, d\tau
\]

\noindent
A system $T$ is \textbf{linear} if $T\{\alpha f + \beta g\} = \alpha\, T\{f\} + \beta\, T\{g\}$.\\
A system $T$ is \textbf{time-invariant} if $T\{f(t-t_0)\} = (Tf)(t-t_0)$.\\
A system that is both is called \textbf{LTI}.

\vspace{14pt}
\noindent\textbf{Prove:}

\begin{enumerate}[leftmargin=2em, itemsep=10pt]

  \item \textbf{Convolution is commutative:}\quad $f \conv g \;=\; g \conv f$

  \item \textbf{Convolution is associative:}\quad $(f \conv g) \conv h \;=\; f \conv (g \conv h)$

  \item \textbf{Convolution is distributive:}\quad $f \conv (g + h) \;=\; f \conv g \;+\; f \conv h$

  \item \textbf{Delta is the identity:}\quad $f \conv \delta \;=\; f$

  \item \textbf{Convolution commutes with differentiation:}
  \[\frac{d}{dt}(f \conv g) \;=\; \frac{df}{dt} \conv g \;=\; f \conv \frac{dg}{dt}\]

  \item \textbf{Support adds:}\quad $\supp(f \conv g) \;\subseteq\; \supp(f) + \supp(g)$

  The \emph{support} of $f$, written $\supp(f)$, is the closure of the set where $f$ is nonzero
  (i.e., where $f$ ``lives,'' including the boundary).
  The $+$ between two sets is the \emph{Minkowski sum}: $A + B = \{a + b : a \in A,\; b \in B\}$.

  In plain words: if $f$ lives on some interval and $g$ lives on some interval,
  show that $f \conv g$ can only be nonzero on an interval whose length is at most
  the sum of those two lengths.

  \textcolor{gray}{\small [Hint: when is the integrand nonzero?]}

  \item \textbf{A system is LTI if and only if it can be expressed by a convolution:}\quad $T$ is LTI $\Leftrightarrow$ $T\{f\} = f \conv h$ for some kernel $h$.

  \textcolor{gray}{\small [Hint: $\Leftarrow$ is easy.\; $\Rightarrow$\, express $f$ using $\delta$]}

  \item \textbf{Gaussians convolve into Gaussians:}
  \[ g_{\sigma_1} \conv g_{\sigma_2} \;=\; g_{\sqrt{\sigma_1^2 + \sigma_2^2}} \,,\quad \text{where}\quad g_\sigma(t) = \tfrac{1}{\sigma\sqrt{2\pi}}\exp\!\left(-\tfrac{t^2}{2\sigma^2}\right) \]

  \textcolor{gray}{\small [Hint: complete the square]}

  \item \textbf{Repeated convolution converges to a Gaussian.}

  In this problem you will show that convolving \emph{any} reasonable filter
  with itself many times produces a Gaussian---a deep connection between
  convolution and the Central Limit Theorem.

  \begin{enumerate}[label=(\alph*), itemsep=6pt]

    \item Consider two \textbf{independent} random variables $X_1, X_2$ with
    PDFs $f_1(x), f_2(x)$.
    Write down the joint PDF $q(x_1, x_2)$ of $(X_1, X_2)$.
    What does independence let you do here?

    \item Let $Y = X_1 + X_2$.
    Express the CDF $F_Y(y) = P(Y \le y)$ as a 2-D integral over the
    region $\{x_1 + x_2 \le y\}$.

    \textcolor{gray}{\small [Hint: set up the limits: $x_1$ from $-\infty$ to $\infty$,
    $x_2$ from $-\infty$ to $y - x_1$.]}

    \item Differentiate $F_Y(y)$ with respect to $y$ to obtain the PDF
    $f_Y(y)$. What familiar operation appears?

    \item Now consider $n$ \textbf{i.i.d.}\ random variables $X_1, \ldots, X_n$,
    each with PDF $f$, mean $\mu$, and finite variance $\sigma^2$.
    Let $S_n = X_1 + \cdots + X_n$.
    Use the result from~(c) to express the PDF of $S_n$ in terms of~$f$.

    \item State the Central Limit Theorem for $S_n$.
    What does it say about $f_{S_n}$ for large~$n$?

    \item Combine (d) and (e): what does this imply about repeatedly
    convolving a nonnegative, normalized function with itself?

    \textcolor{gray}{\small [This is the punchline: the $n$-fold self-convolution
    of any such function converges to a Gaussian, with variance growing
    linearly in~$n$.]}

    \item \textbf{(Optional---numerical verification.)}
    Take the box filter $h = [\tfrac{1}{3},\, \tfrac{1}{3},\, \tfrac{1}{3}]$
    and convolve it with itself 30 times numerically.
    Plot the result and overlay a Gaussian with the predicted variance.
    Does it match?

  \end{enumerate}

\end{enumerate}

\vfill
\begin{center}
{\small\textcolor{gray}{Modern Computer Vision --- Technion --- Spring 2026 --- Assaf Shocher}}
\end{center}

\end{document}
